% paper/sections/12f_pressure_solver.tex
%
% Tier V: 時間積分の検証
%   U8 = Time integration suite
%        (a) TVD-RK3 ODE 3 次 + 安定領域
%        (b) AB2 ODE 2 次
%        (c) CN + ADI predictor 拡散 MMS 2 次
%        (d) 粘性 3-layer Layer A/B/C 個別精度
% ═══════════════════════════════════════════════════════════════════════════════
\subsection{Tier V：時間積分}
\label{sec:verify_time}

NS パイプラインの時間進展は，
CLS 移流の TVD-RK3（§~\ref{sec:time_int}），
NS Predictor の IMEX-BDF2 / AB2（§~\ref{sec:time_level1}），
粘性項の Crank--Nicolson + ADI sweep（§~\ref{sec:time_int}）の
3 系統に分かれる．
本 Tier の単体テスト U8 はこれらを 4 sub-test で suite 検証する．

% ─────────────────────────────────────────────────────────────────────────────
\subsubsection{U8：Time integration suite}
\label{sec:U8_time_integration_suite}
% ─────────────────────────────────────────────────────────────────────────────

\noindent\textbf{目的：}
TVD-RK3 / AB2 / CN+ADI predictor / 粘性 3-layer の 4 sub-test で
時間積分スキームの設計精度を ODE / MMS テストで検証する．

\noindent\textbf{Part 2 対応 primitive：}
\#13 (TVD-RK3),\#14 (IMEX-BDF2),\#15 (ADI predictor sweep),\#11 (粘性 3-layer)．

\noindent\textbf{下流連関：}
全 NS time 系（KE decay,TGV temporal）．

% ───────────────
\paragraph{(a) TVD-RK3 ODE 3 次精度＋安定領域}
\label{sec:verify_tvd_rk3}

スカラー ODE $\mathrm{d}q/\mathrm{d}t = -q$，$q(0) = 1$，厳密解 $q(1) = e^{-1}$；
ステップ数 $n = 4, 8, \ldots, 512$．

\begin{table}[htbp]
\centering
\caption{TVD-RK3 の時間精度}
\label{tab:tvd_rk3}
\renewcommand{\arraystretch}{1.2}
\footnotesize
\begin{tabular}{rc cc}
\toprule
$n$ & $\Delta t$ & $|q - e^{-1}|$ & 次数 \\
\midrule
  4 & $2.50 \times 10^{-1}$ & $2.93 \times 10^{-4}$  & ---    \\
  8 & $1.25 \times 10^{-1}$ & $3.31 \times 10^{-5}$  & $3.14$ \\
 16 & $6.25 \times 10^{-2}$ & $3.93 \times 10^{-6}$  & $3.07$ \\
 32 & $3.13 \times 10^{-2}$ & $4.80 \times 10^{-7}$  & $3.04$ \\
 64 & $1.56 \times 10^{-2}$ & $5.92 \times 10^{-8}$  & $3.02$ \\
128 & $7.81 \times 10^{-3}$ & $7.35 \times 10^{-9}$  & $3.01$ \\
256 & $3.91 \times 10^{-3}$ & $9.17 \times 10^{-10}$ & $3.00$ \\
512 & $1.95 \times 10^{-3}$ & $1.14 \times 10^{-10}$ & $3.00$ \\
\bottomrule
\end{tabular}
\end{table}

収束次数は漸近的に $3.00$ に近づき，設計通りの $\Ord{\Delta t^3}$ を達成．
固定格子移流テスト（CCD 周期 $N = 256$）では空間離散のスペクトル半径による
安定領域制約を確認した．

\noindent\textbf{結論：}$\Ord{\Delta t^{3.00}}$．\textbf{合格}．

% ───────────────
\paragraph{(b) AB2 ODE 2 次精度}
\label{sec:verify_ab2}

スカラー ODE $\mathrm{d}q/\mathrm{d}t = -q$，AB2（前進 Euler 起動）；
$n = 16, 32, \ldots, 512$．

\begin{table}[htbp]
\centering
\caption{AB2 の時間精度}
\label{tab:ab2_time}
\renewcommand{\arraystretch}{1.2}
\footnotesize
\begin{tabular}{rc cc}
\toprule
$n$ & $\Delta t$ & $|q - e^{-1}|$ & 次数 \\
\midrule
 16 & $6.25 \times 10^{-2}$ & $1.42 \times 10^{-4}$ & ---    \\
 32 & $3.13 \times 10^{-2}$ & $3.27 \times 10^{-5}$ & $2.12$ \\
 64 & $1.56 \times 10^{-2}$ & $7.84 \times 10^{-6}$ & $2.06$ \\
128 & $7.81 \times 10^{-3}$ & $1.92 \times 10^{-6}$ & $2.03$ \\
256 & $3.91 \times 10^{-3}$ & $4.73 \times 10^{-7}$ & $2.02$ \\
512 & $1.95 \times 10^{-3}$ & $1.18 \times 10^{-7}$ & $2.01$ \\
\bottomrule
\end{tabular}
\end{table}

収束次数は $\Ord{\Delta t^2}$ に漸近．寄生根 $|\rho_2| = 0.5$ は安定領域内で，
前進 Euler 起動の $\Ord{\Delta t}$ 誤差は数ステップで減衰．
EXT2 外挿（IMEX-BDF2 の対流処理）も同じ $(3/2, -1/2)$ 係数の 2 次零安定多段法．

\noindent\textbf{結論：}$\Ord{\Delta t^{2.0}}$．\textbf{合格}．

% ───────────────
\paragraph{(c) Crank--Nicolson + ADI predictor sweep 拡散 MMS 2 次}
\label{sec:verify_cn_viscous}

2D 拡散方程式 $u_t = \nu \nabla^2 u$（$\nu = 0.01$，周期境界）に対し
厳密解 $u = \exp(-8\pi^2 \nu t)\sin(2\pi x)\sin(2\pi y)$ で $N = 64$ 固定，
$\Delta t = T/10 \to T/320$ で時間収束を計測．

\noindent\textbf{結果：}
収束次数は $T/10 \to T/160$ で $2.00 \to 1.99 \to 2.04 \to 2.18$ と
安定して $\Ord{\Delta t^2}$ を示し，
$T/320$ で空間離散の丸め誤差床に到達．

\noindent\textbf{無条件安定性：}
$\mathrm{CFL}_\nu = \nu\Delta t/h^2 = 0.5$--$10$ で
CN は $\mathrm{CFL}_\nu = 10$ でも単調減少を維持，
前進 Euler は $\mathrm{CFL}_\nu = 5$ で初期値の約 $600$ 倍に発散する．
ADI 方向分裂誤差は第~\ref{sec:verification}章 の TGV 結合テストで検証．

\noindent\textbf{結論：}CN $\Ord{\Delta t^{2.0}}$＋無条件安定．\textbf{合格}．

% ───────────────
\paragraph{(d) 粘性 3-layer Layer A/B/C 個別精度（交差微分粘性項）}
\label{sec:verify_cross_viscous}

CN は対角粘性項 $\partial_x[\mu\,\partial_x u]$ のみ暗黙処理し，
交差微分項 $\partial_x[\mu\,\partial_y u]$ は陽的評価する（§~\ref{sec:time_int}）．
粘性比 $\mu_l/\mu_g$ がこの分離に与える影響を 3 層に分けて MMS で測定する：
Layer A（対角項のみ），Layer B（対角+陽的交差項），Layer C（界面 HFE 統合）．

\noindent\textbf{結果：}
\begin{itemize}[noitemsep]
  \item Layer A: $\mu_l/\mu_g = 1$ で LTE $\Ord{\Delta t^2}$（古典 CN）
  \item Layer B: $\mu_l/\mu_g = 10, 100$ で LTE 傾き $1.00$ に低下（$\Ord{\Delta t}$ 劣化）
  \item Layer C: 1D MMS proxy では Layer B と同じ $\Ord{\Delta t}$
        （cross-EE 構造由来）；
        2D ADI 実装（第~\ref{sec:verification}章 TGV 結合テスト）では
        HFE 平滑化 $\times$ ADI 分裂誤差の相互作用で
        $\Ord{\Delta t^{1.5}}$ 程度回復
\end{itemize}
全体の粘性ソルバ精度は
$\Ord{\Delta t^2}$（対角 CN）＋ $\Ord{\mu_l/\mu_g \cdot \Delta t}$（交差項）．
$\mu_l/\mu_g \lesssim 5$ では交差項寄与は CSF $\Ord{h^{1\text{-}2}}$ と同程度であり
実用上支配的ではない．

\noindent\textbf{結論：}
高粘性比界面近傍で $\Ord{\Delta t}$ 劣化を確認；
理論予測（§~\ref{sec:time_int}）と一致．
\textbf{合格}（理論一致）／実用的に許容可能．
